\documentclass[11pt,a4paper]{article}
\usepackage[utf8]{inputenc}
\usepackage[T1]{fontenc}
\usepackage[brazil]{babel}
\usepackage{lmodern}
\usepackage{microtype}
\usepackage[margin=2.3cm]{geometry}
\usepackage{booktabs,longtable,tabularx,array}
\usepackage{xcolor}
\usepackage{tikz}
\usetikzlibrary{positioning,arrows.meta,fit,backgrounds}
\usepackage{enumitem}
\usepackage{caption}
\usepackage{fancyhdr}
\usepackage{amssymb}
\usepackage[hidelinks]{hyperref}
\emergencystretch=3em

\definecolor{latao}{HTML}{B9802A}
\definecolor{tinta}{HTML}{0B1220}
\definecolor{lagoa}{HTML}{1F8A80}
\definecolor{coral}{HTML}{C9533A}
\definecolor{cinza}{HTML}{5B6676}
\definecolor{violeta}{HTML}{7B61D1}

\pagestyle{fancy}\fancyhf{}
\lhead{\small\color{cinza}HRAI-GO · RT-VIS-03 · documento de estado-alvo}
\rhead{\small\color{cinza}\thepage}
\renewcommand{\headrulewidth}{0.3pt}
\newcommand{\nota}[1]{\par\smallskip\noindent\fcolorbox{latao}{latao!6}{\parbox{\dimexpr\linewidth-2\fboxsep-2\fboxrule}{\small #1}}\par\smallskip}
\newcommand{\meta}[1]{{\color{violeta}\textbf{[meta: #1]}}}
\newcommand{\rf}[1]{{\color{cinza}\footnotesize(#1)}}
\newcolumntype{L}{>{\raggedright\arraybackslash}X}

\title{\textbf{HRAI-GO: Human-centred Regulatory AI Governance Observatory}\\[5pt]
\large Descrição metodológica e aplicação do observatório\\[2pt]
\normalsize Relatório técnico escrito no presente, como estado-alvo verificável}
\author{Plataforma HRAI-GO · Aeonbridge}
\date{RT-VIS-03 · versão de 17 de novembro de 2026 · referência de aceite da ESP-ONT-02}

\begin{document}
\maketitle
\thispagestyle{empty}

\nota{\textbf{Como ler este documento.} Ele descreve o HRAI-GO como a plataforma deve existir quando a especificação ESP-ONT-02 estiver implementada. Está escrito no presente de propósito (técnica de ``escrever o futuro como queremos'').
\begin{itemize}[leftmargin=*]
  \item \textbf{É uma referência de aceite, não um relato de resultados.} Cada afirmação sobre o que a plataforma faz corresponde a um requisito, indicado entre parênteses, e a uma verificação, listada no checklist do apêndice~\ref{ap:checklist}.
  \item \textbf{Números em \meta{\ldots} são metas}, não medições.
  \item \textbf{Únicos números medidos:} os resultados que já existem no repositório em 17/11/2026, como a classificação das teses e as contagens do RT-ONT-01. Eles estão identificados como tais.
\end{itemize}
A plataforma está pronta quando todos os itens do checklist estiverem marcados.}

\begin{abstract}
\noindent O HRAI-GO é um observatório público de governança de inteligência artificial. Ele representa os modelos de regulação de IA do Brasil e de outras jurisdições numa ontologia com proveniência até a fonte e os compara por uma taxonomia publicada e revisada por pares, estendida e declarada. Raciocina sobre eles com regras auditáveis que nunca emitem veredicto jurídico e põe o resultado diante de pessoas, ao vivo, num processo de deliberação em que perspectivas sintéticas ampliam o que é considerado sem nunca decidir. Cada elemento que o observatório mostra pode ser clicado e responde ``o que é, de onde veio e por que está aqui''. Este relatório descreve o método e a aplicação, e serve como critério de aceite da plataforma.
\end{abstract}

\tableofcontents
\newpage

%=====================================================================
\section{O nome e o que ele promete}

O nome é uma decisão de projeto. Cada parte corresponde a um compromisso verificável.

\begin{center}\small
\begin{tabularx}{\linewidth}{c >{\raggedright\arraybackslash}p{36mm} L}
\toprule
Letra & Significado & Por quê, e onde isso aparece na plataforma\\
\midrule
\textbf{H} & \textbf{Human-centred} & O humano vem primeiro: supervisão humana, papéis humanos que contestam, qualificam ou apoiam com evidência, e revisão humana antes de publicar. O termo é alinhado ao AI Act, à OCDE e ao PL~2338. Na plataforma: nada extraído vira fato sem decisão de revisão humana; nenhuma regra emite veredicto; toda deliberação termina em decisão humana registrada.\\
\textbf{R} & \textbf{Regulatory} & Inteligência regulatória: normas, versões, dispositivos, obrigações, tensões e lacunas. Na plataforma: o núcleo ontológico representa cada instrumento em suas versões, com cada dispositivo ligado à captura da fonte.\\
\textbf{AI} & \textbf{AI} & A IA é ao mesmo tempo o objeto regulado (aplicações de IA) e o meio de análise (agentes e Shadow Council). Na plataforma: aplicações aparecem como alvos de propagação de impacto, e agentes geram extrações candidatas e perspectivas sintéticas sempre rotuladas.\\
\textbf{GO} & \textbf{Governance Observatory} & É o papel público da plataforma: observar, comparar e expor a governança de IA, do Brasil e do mundo, de forma aberta. Na plataforma: o site é público, os identificadores resolvem, as fontes estão disponíveis e cada comparação mostra o escopo em que vale.\\
\bottomrule
\end{tabularx}
\end{center}

\noindent O nome completo, \emph{HRAI-GO: Human-centred Regulatory AI Governance Observatory}, aparece no cabeçalho do site, nos metadados da página, no README do repositório e nas ontologias \rf{V-01}.

%=====================================================================
\section{O problema que o observatório resolve}

A governança de IA virou um campo de comparações imprecisas. ``Regulação de IA'' designa ao mesmo tempo leis vinculantes, planos de investimento, princípios não estatutários e guias voluntários. Análises influentes circulam com teses fortes e pouca evidência rastreável. Decisões coletivas sobre o tema costumam ser tomadas em salas onde os grupos mais afetados não estão.

O HRAI-GO responde a isso com três garantias:
\begin{enumerate}[leftmargin=*]
  \item \textbf{Comparar o comparável.} Cada instrumento é classificado pelo que é (lei, proposta, plano, princípio, guia) e por como regula, com evidência por dispositivo.
  \item \textbf{Separar fato, inferência e opinião.} Fatos têm fonte, inferências têm regra e cadeia, opiniões entram como afirmações com evidência e contra-evidência.
  \item \textbf{Pôr o raciocínio diante de pessoas.} Quem decide é humano. A máquina amplia o que é considerado e mostra onde está o desacordo.
\end{enumerate}

%=====================================================================
\section{Arquitetura}

\begin{figure}[ht]
\centering
\begin{tikzpicture}[node distance=6mm and 8mm, every node/.style={font=\footnotesize},
  mod/.style={draw=tinta!70, rounded corners=4pt, align=center, text width=40mm, minimum height=24mm, fill=white},
  hum/.style={draw=latao, rounded corners=4pt, fill=latao!10, align=center, text width=40mm, minimum height=12mm},
  seta/.style={-{Stealth[length=2mm]}, thick, color=tinta!70}]
\node[mod] (nucleo) {\textbf{Núcleo regulatório}\\jurisdições, instrumentos, versões, dispositivos, requisitos, status, contexto, proveniência, afirmações};
\node[mod, right=of nucleo] (comp) {\textbf{Módulo comparativo}\\taxonomia de Alanoca et al.\ + E1--E7, codificações, correspondências, áreas substantivas};
\node[mod, right=of comp] (scse) {\textbf{SCSE}\\tensões de deliberação, conselho, estratos, assentos, rodadas, posições, decisões};
\node[mod, below=10mm of comp, text width=60mm, minimum height=12mm] (raz) {\textbf{Raciocínio rastreável}\\regras R01--R30 + consultas PC1--PC8, com justificativa até a fonte};
\node[hum, below=8mm of raz, text width=60mm] (ui) {\textbf{Interfaces}: câmara ao vivo, celular do participante, comparador, monitoramento, base de conhecimento};
\node[hum, left=of ui, text width=34mm] (rev) {\textbf{Revisão humana}\\candidato $\to$ aprovado $\to$ publicado};
\node[hum, right=of ui, text width=34mm] (dec) {\textbf{Decisão humana}\\registro com justificativa};
\draw[seta] (nucleo) -- (comp); \draw[seta] (scse) -- (comp);
\draw[seta] (nucleo.south) |- (raz.west); \draw[seta] (comp) -- (raz); \draw[seta] (scse.south) |- (raz.east);
\draw[seta] (raz) -- (ui); \draw[seta] (rev) -- (ui); \draw[seta] (ui) -- (dec);
\end{tikzpicture}
\caption{Arquitetura do HRAI-GO. Em latão, os pontos em que só humanos atuam.}
\end{figure}

\begin{itemize}[leftmargin=*]
  \item \textbf{Núcleo regulatório (HRAI-GO, v1.2.0).} Representa o que as fontes dizem. Cada instrumento existe em versões, cada versão contém dispositivos, e cada dispositivo deriva de uma captura com hash SHA-256. Normas vigentes, propostas e planos são classes disjuntas, de modo que uma proposta nunca se apresenta como obrigação \rf{RF-03}.
  \item \textbf{Módulo comparativo.} Classifica os instrumentos para comparação, usando a taxonomia de Alanoca et al.\ (2025) com as extensões declaradas (seção~\ref{sec:taxonomia}) \rf{RF-16}.
  \item \textbf{SCSE (Shadow Council Simulation Environment).} Modela a deliberação: tensões que a sala discute, assentos sintéticos por estrato, cadeiras vazias, rodadas, posições humanas e sintéticas, e decisões. Importa o núcleo em vez de redeclarar seus conceitos \rf{RF-17}.
  \item \textbf{Fonte única.} Tudo o que as interfaces mostram é gerado do grafo no build. Não existe cadastro paralelo editado à mão. O contador de elementos sem lastro está em zero \rf{RF-09, RF-12}.
\end{itemize}

%=====================================================================
\section{Método}

\subsection{M1 · Corpus com proveniência}
O corpus fica em \texttt{docs/ai-regulations/} e cobre:
\begin{itemize}[leftmargin=*]
  \item \textbf{Brasil:} o PBIA 2024--2028; o PL~2338/2023 no texto inicial e no substitutivo aprovado pelo Senado em 10/12/2024, como duas versões do mesmo instrumento \rf{RF-01}.
  \item \textbf{União Europeia:} o AI Act.
  \item \textbf{Reino Unido:} o white paper e a análise de impacto.
  \item \textbf{EUA:} o NIST AI RMF.
  \item \textbf{China:} as jurisdições adicionais decididas no RF-02. Nenhuma jurisdição aparece no observatório sem ao menos uma captura com hash \rf{RF-02}.
\end{itemize}
Cada documento entra como \texttt{SourceSnapshot}, com URI, data de captura e hash. Quando a fonte muda, nasce uma nova captura; a anterior não é apagada.

\subsection{M2 · Extração candidata e revisão humana}
Agentes de extração propõem dispositivos, requisitos, codificações, correspondências e indicadores. Tudo entra como \texttt{CandidateAssertion} num lote de extração. Um revisor humano aprova ou rejeita com \texttt{ReviewDecision}. A publicação num \texttt{KnowledgeGraphRelease} é um ato explícito e autorizado, nunca uma inferência. Agentes não aprovam o que extraíram \rf{RF-04, invariante I-3}.

\subsection{M3 · Taxonomia de base e extensões}\label{sec:taxonomia}
\nota{\textbf{Declaração de método.} Para a camada comparativa, o HRAI-GO utiliza, onde possível, a taxonomia de Alanoca, Gur-Arieh, Zick e Klyman, \emph{Comparing Apples to Oranges: A Taxonomy for Navigating the Global Landscape of AI Regulation} (FAccT~'25, doi:10.1145/3715275.3732059).
\begin{itemize}[leftmargin=*]
  \item \textbf{Por que esta taxonomia:} tem o mesmo objetivo, comparar regulações de IA entre jurisdições; usa critérios mensuráveis e comuns aos documentos; foi revisada por pares e validada com reguladores e especialistas; e já foi aplicada ao AI Act e ao PL~2338.
  \item \textbf{Onde foi usada:} os 25 critérios, em 12 categorias, são usados com as definições originais.
  \item \textbf{Por que foi estendida:} para construir a ontologia foi necessário estendê-la em sete pontos, que comportam os casos abaixo.
\end{itemize}}

\begin{center}\small
\begin{tabularx}{\linewidth}{l >{\raggedright\arraybackslash}p{48mm} L}
\toprule
& Extensão & Caso comportado\\
\midrule
E1 & Forma jurídica, incluindo instrumentos não vinculantes & PBIA (plano), Reino Unido (princípios não estatutários), NIST (guia voluntário)\\
E2 & Objetivo primário & o caráter desenvolvimentista e de soberania do modelo brasileiro\\
E3 & Cobertura substantiva (28 áreas + IA generativa/GPAI) & lacunas de conteúdo: energia, crianças, desigualdade regional, participação\\
E4 & Codificação em escala, com evidência por valor & valores comparáveis e rastreáveis, em vez de descrições\\
E5 & Codificação por versão do instrumento & PL~2338 de 2023 $\times$ 2024; instrumentos revogados\\
E6 & Capacidade de implementação & orçamento e pessoal de supervisão (Eixo~5 do PBIA)\\
E7 & Participação na implementação & deliberação pública ao vivo e Shadow Council\\
\bottomrule
\end{tabularx}
\end{center}

Cada valor é um \texttt{CriterionCoding}: versão do instrumento, critério, valor na escala aprovada, dispositivo que justifica, codificador e revisão. Para o Brasil e a UE, as codificações dos autores do paper estão registradas ao lado das nossas. A concordância é publicada por critério \meta{concordância publicada para os 25 critérios nas duas jurisdições}, e cada divergência foi resolvida contra o texto primário \rf{RF-16}.

\subsection{M4 · Correspondência entre dispositivos}
Para cada área substantiva, os dispositivos do Brasil e do AI Act estão ligados por \texttt{ProvisionCorrespondence}, de um destes tipos: \emph{equivalente}, \emph{adaptado}, \emph{parcial} ou \emph{divergente}. Cada correspondência tem evidência e revisão humana. É sobre essas correspondências que o observatório responde ``o que o Brasil incorporou'' \rf{RF-06}.

\subsection{M5 · Raciocínio rastreável}
O motor executa as 30 regras SWRL (R01--R30) lidas da própria ontologia, até o ponto fixo, com fechamento explícito de subclasse, subpropriedade e inversa. Cada tripla inferida guarda a regra, as ligações e as triplas do corpo, e a pergunta ``por quê?'' devolve a árvore até as triplas declaradas e a captura da fonte. O raciocínio do grafo completo roda no CI em menos de 5\,s \rf{RNF-02}.

Invariantes garantidas por teste:
\begin{itemize}[leftmargin=*]
  \item \textbf{I-1:} nenhuma regra afirma status jurídico, conformidade, publicação ou decisão;
  \item \textbf{I-2:} cenários e afirmações nunca se misturam com instâncias de fonte;
  \item \textbf{I-3:} nenhum agente cria decisão de revisão;
  \item \textbf{I-4:} sintético nunca é contado como humano.
\end{itemize}

\subsection{M6 · Perguntas de competência}
O observatório responde a oito perguntas. Cada uma tem uma consulta versionada, um teste de regressão no CI e uma tela \rf{RF-08}:

\begin{center}\small
\begin{tabularx}{\linewidth}{l L L}
\toprule
& Pergunta & Onde a resposta aparece\\
\midrule
PC1 & Que áreas cada modelo cobre, e com que força jurídica? & Comparador: matriz jurisdição $\times$ área\\
PC2 & Que dispositivos do Brasil correspondem, adaptam ou divergem do AI Act? & Comparador: lado a lado\\
PC3 & O que o Brasil incorporou que pressupõe capacidade que o contexto não tem? & Câmara: ``incorporado sem capacidade''\\
PC4 & O que o contexto expõe a risco e nenhum instrumento cobre? & Câmara: lacunas de exposição, com escopo\\
PC5 & Como cada modelo se classifica nos critérios da taxonomia, e com que evidência? & Ficha do modelo; comparador\\
PC6 & Que afirmações sustentam cada recomendação, quem as contesta, com que evidência? & Mural de tensões; fichas\\
PC7 & Quando um texto muda, o que é reaberto? & Monitoramento\\
PC8 & Que posições humanas e sintéticas apoiam ou contestam cada alternativa, e o que mudou na discussão? & Mesa da sala; relatório da sessão\\
\bottomrule
\end{tabularx}
\end{center}

\subsection{M7 · Afirmações com contra-evidência}
Análises externas não entram como fato. Cada tese vira \texttt{PolicyClaim}, ligada às evidências do corpus que a sustentam e às contra-afirmações que a contestam, com relações distintas para contestar, qualificar e complementar. O observatório mostra o estado de cada tese: \emph{contestada}, \emph{com evidência} ou \emph{sem evidência no corpus}. Nunca diz ``verdadeira'' ou ``falsa'' \rf{RF-11}.

\nota{\textbf{Resultado já medido (17/11/2026).} As 12 teses do documento ``Análise Comparativa e Diretrizes Estratégicas'' foram processadas pelo método: 4 contestadas, 6 com evidência e 2 sem evidência no corpus (ESP-ONT-02, apêndice~A). No estado-alvo, essa classificação é recalculada automaticamente a cada versão do corpus. Por exemplo, a tese C04 é reavaliada quando o substitutivo do Senado entra como fonte primária.}

\subsection{M8 · Deliberação humana com o Shadow Council}
O Shadow Council é um conjunto de 20 perspectivas sintéticas em 8 estratos. Cada perspectiva é ancorada em documentos e declara seus pontos cegos, e os grupos sem fonte aparecem como cadeiras vazias. Ele participa de um protocolo de cinco rodadas:

\begin{enumerate}[leftmargin=*, start=0]
  \item \textbf{Enquadramento:} a tensão da sessão vem do grafo, com afirmações e evidências.
  \item \textbf{Posição cega:} cada pessoa se posiciona antes de ver o conselho.
  \item \textbf{Revelação:} o conselho mostra as posições por estrato, com fontes e cadeiras vazias, e destaca os estratos ausentes da sala.
  \item \textbf{Revisão:} as pessoas apoiam, contestam, pedem fonte e podem mudar de posição.
  \item \textbf{Decisão:} registro humano que cita as afirmações em que se apoia e preserva as divergências.
\end{enumerate}

O conselho não vota, não recomenda e não produz ``consenso sintético''. Cada sessão gera um relatório com quatro métricas \rf{RF-10}:
\begin{itemize}[leftmargin=*]
  \item \textbf{amplitude:} afirmações consideradas;
  \item \textbf{ancoragem:} fração das afirmações citadas que têm evidência;
  \item \textbf{revisão de posição:} quantas posições mudaram entre as rodadas 1 e 3;
  \item \textbf{divergência preservada:} quantas alternativas terminam com desacordo registrado.
\end{itemize}

\subsection{M9 · Monitoramento contínuo}
Cada nova captura gera um \texttt{RegulatoryChangeEvent}. As regras R10 a R13 propagam o impacto \emph{potencial} até dispositivos, requisitos, avaliações e aplicações. O monitoramento lista o que foi reaberto e alimenta a fila de revisão humana. Codificações da taxonomia ligadas à versão anterior são marcadas para recodificação \rf{RF-01, E5}.

%=====================================================================
\section{Aplicação}

\subsection{Uma sessão pública ao vivo}
Um auditório com representantes de governo, reguladores, academia, sociedade civil e setor privado. No telão, a câmara 3D, aberta com \texttt{?telao=1}. O QR code leva a \texttt{go.aeonbridge.io/ai-ethic-br}.

\begin{itemize}[leftmargin=*]
  \item \textbf{Entrada.} Quem escaneia escolhe o setor e entra na câmara. Recebe um avatar gerado a partir do identificador anônimo, igual em todos os aparelhos e sem nenhum dado pessoal, e um apelido como ``Academia·7f3''. A luz da pessoa aparece no anel externo, no arco do seu setor.
  \item \textbf{Participação pelo celular.} No celular, a pessoa é participante, não espectadora.
    \begin{itemize}
      \item \emph{Minhas escolhas:} escolhe a tensão, se posiciona no eixo, aponta a lacuna e publica uma contribuição.
      \item \emph{Mesa:} lê as contribuições dos outros e reage com apoio, contestação ou pedido de fonte.
      \item \emph{Agora:} acompanha o que a sala faz.
    \end{itemize}
    Tudo funciona com o polegar, com alvos de toque de pelo menos 44\,px, e cada ação aparece para todos e no telão.
  \item \textbf{Fichas com ``por quê?''.} Tocar em qualquer elemento abre a ficha dele: um bloco da torre, um assento, um arco de tensão, um anel da cúpula, uma pessoa. A ficha mostra o identificador na ontologia, a fonte com hash, o estado de revisão e, para inferências, a cadeia até a fonte. Nenhum elemento exibido está ``sem lastro'' \rf{RF-12}.
  \item \textbf{Deliberação.} O apresentador conduz as cinco rodadas. Na revelação, o conselho mostra, por exemplo, que nenhum participante da sala representa ``grupos diretamente afetados'' e que ``povos indígenas'' é cadeira vazia por falta de fonte. A decisão final é registrada pela coordenação humana e entra no grafo como candidata.
\end{itemize}

\subsection{O comparador Brasil $\times$ AI Act}
Um analista abre o comparador e vê:
\begin{itemize}[leftmargin=*]
  \item a matriz dos 25 critérios para o Brasil (substitutivo de 2024) e para a UE, cada célula com o dispositivo de origem e a concordância com a codificação dos autores do paper \rf{PC5};
  \item para cada área, os dispositivos lado a lado, com o tipo de correspondência \rf{PC2};
  \item a lista de requisitos transpostos que pressupõem capacidade sem indicador suficiente, por exemplo o orçamento de supervisão frente ao Eixo~5 do PBIA, sempre com a fonte do indicador e o limiar \rf{PC3};
  \item as exposições do contexto sem instrumento que as cubra, dentro do escopo declarado da avaliação de cobertura \rf{PC4};
  \item o que é próprio do Brasil: as áreas e os objetivos que só o modelo brasileiro tem, como a soberania computacional com orçamento, a nuvem de governo e o LLM em português (E2 e E3).
\end{itemize}

\subsection{Quando o texto muda}
O substitutivo do PL~2338 chega ao corpus como nova versão. O monitoramento mostra o evento, os dispositivos potencialmente afetados, as codificações que precisam ser refeitas e as teses cuja classificação mudou. A tese C04 é reavaliada contra o texto primário que cria o comitê de especialistas. Nada é atualizado como fato sem revisão humana \rf{PC7}.

\subsection{A base de conhecimento}
A base de conhecimento expõe os mesmos dados para leitura: instrumentos por jurisdição, linha do tempo, componentes do modelo brasileiro com força jurídica e fonte, e as teses com seu estado. Os identificadores das ontologias e instâncias resolvem no navegador \rf{RNF-01}.

%=====================================================================
\section{Como o ``H'' é garantido}
\begin{center}\small
\begin{tabularx}{\linewidth}{>{\raggedright\arraybackslash}p{48mm} L}
\toprule
Compromisso \emph{human-centred} & Mecanismo verificável\\
\midrule
Revisão humana antes de publicar & \texttt{CandidateAssertion} $\to$ \texttt{ReviewDecision} humana $\to$ release explícito; invariante I-3 testada\\
Nenhum veredicto automático & invariante I-1 testada em todas as regras; ressalva em cada regra\\
Humanos contestam, qualificam e apoiam com evidência & reações da Mesa gravadas como posições sobre afirmações; relações de contestar, qualificar e complementar\\
Sintético nunca passa por humano & disjunção na TBox, regra R21, cor e rótulo distintos na interface; invariante I-4\\
Decisão é humana e justificada & \texttt{DecisionRecord} com afirmações citadas e divergências preservadas\\
Ausências à vista & cadeiras vazias e estratos ausentes da sala destacados na revelação\\
Privacidade & participação anônima, sem nome nem e-mail; avatar derivado do identificador anônimo\\
\bottomrule
\end{tabularx}
\end{center}

%=====================================================================
\section{Limites declarados}
Mesmo implementado por inteiro, o HRAI-GO tem limites, declarados na própria interface:
\begin{itemize}[leftmargin=*]
  \item \textbf{Não é aconselhamento jurídico.} Marcar algo como ``potencialmente afetado'' não diz se há conformidade ou violação.
  \item \textbf{Perspectivas sintéticas não são representativas da população.} Elas ampliam o que é considerado; não substituem consulta pública.
  \item \textbf{Lacunas valem dentro do escopo declarado.} Uma ausência só vale dentro do escopo da avaliação de cobertura que a produziu.
  \item \textbf{Validade das codificações.} As codificações dependem das escalas aprovadas e da revisão humana. A concordância com codificadores externos é publicada justamente para tornar isso visível.
\end{itemize}

%=====================================================================
\appendix
\section{Checklist de aceitação}\label{ap:checklist}
\noindent Cada linha é uma afirmação deste documento. A plataforma está conforme quando todas estiverem marcadas com evidência anexada (link do CI, captura de tela ou commit).

\begin{center}\footnotesize
\begin{longtable}{l >{\raggedright\arraybackslash}p{62mm} >{\raggedright\arraybackslash}p{70mm} c}
\toprule
Id & Afirmação & Como verificar & OK\\
\midrule\endhead
V-01 & O nome HRAI-GO: Human-centred Regulatory AI Governance Observatory aparece no site, no README e nas ontologias & inspecionar cabeçalho e metadados do site, README e \texttt{dcterms:title} das ontologias & $\square$\\
V-02 & Corpus com BR (PBIA, PL 2023 e 2024), UE, Reino Unido e EUA, cada documento com captura e hash & listar \texttt{SourceSnapshot} no grafo; conferir SHA-256 dos arquivos & $\square$\\
V-03 & Decisão sobre a China registrada; nenhuma jurisdição sem fonte & consulta: jurisdições exibidas sem snapshot = 0 & $\square$\\
V-04 & Três módulos sem classes duplicadas; SCSE importa o núcleo & teste RF-17 no CI & $\square$\\
V-05 & TBox v1.2.0 com R24--R30; invariantes I-1 a I-4 testadas & \texttt{tests/verify\_regulatory\_ontology.py} verde & $\square$\\
V-06 & 25 critérios de Alanoca et al.\ com definições originais + E1--E7 declaradas, com crédito e DOI & inspecionar esquema \texttt{RegulatoryDesignCriterion} & $\square$\\
V-07 & Escalas da E4 aprovadas por humano & \texttt{ReviewDecision} humana para cada escala & $\square$\\
V-08 & BR e UE codificados nos 25 critérios, cada valor com dispositivo de origem & consulta: codificações sem evidência = 0 & $\square$\\
V-09 & Concordância com os autores publicada por critério; divergências resolvidas contra o texto primário & página do comparador; registro das resoluções & $\square$\\
V-10 & Correspondências BR $\leftrightarrow$ UE revisadas para as áreas cobertas pelas duas & PC2 responde; correspondências sem revisão = 0 & $\square$\\
V-11 & Indicadores de capacidade e exposição com fonte e ano-base & PC3 e PC4 respondem com resultado ou ``sem indicador'' & $\square$\\
V-12 & PC1 a PC8 versionadas e testadas no CI & CI verde; teste falha se uma PC deixar de responder & $\square$\\
V-13 & Raciocínio completo em menos de 5\,s & tempo medido no CI & $\square$\\
V-14 & ``Por quê?'' em toda inferência exibida, até a captura da fonte & abrir 5 fichas de inferência ao acaso & $\square$\\
V-15 & Interface gerada do grafo; contador ``sem lastro'' = 0 & build publica o contador; \texttt{modelo.js} não editado à mão & $\square$\\
V-16 & Tensões, estratos, assentos e cenários na SCSE, ligados ao núcleo & consulta: sementes sem ligação = 0 & $\square$\\
V-17 & Teses externas como afirmações com estado recalculado a cada versão & \texttt{scripts/ontologia/afirmacoes.py} no CI & $\square$\\
V-18 & Celular participante: escolhas, Mesa, Agora; alvos $\geq$ 44\,px; ações visíveis para todos e no telão & teste em aparelho real + teste automatizado de toque & $\square$\\
V-19 & Fichas em todos os elementos visíveis da câmara & tocar torre, assento, arco, cúpula, pessoa, mesa & $\square$\\
V-20 & Protocolo de cinco rodadas, com posição cega antes do conselho & ensaio completo de uma sessão & $\square$\\
V-21 & Relatório de sessão com as quatro métricas & gerado ao final do ensaio & $\square$\\
V-22 & Decisão humana registrada com afirmações citadas e divergências preservadas & \texttt{DecisionRecord} do ensaio & $\square$\\
V-23 & Nova versão do PL reabre dispositivos, codificações e teses & inserir a versão e conferir o monitoramento & $\square$\\
V-24 & Identificadores resolvem no navegador & abrir 3 IRIs ao acaso & $\square$\\
V-25 & Limites declarados visíveis na interface & inspecionar telas do comparador e das fichas & $\square$\\
\bottomrule
\end{longtable}
\end{center}

\section{Rastreabilidade}
\begin{itemize}[leftmargin=*]
  \item \textbf{Especificação:} \texttt{docs/especificacoes/ESP-ONT-02/especificacao.pdf} (requisitos RF-01 a RF-17, RNF-01 e RNF-02).
  \item \textbf{Roteiro de execução:} \texttt{docs/especificacoes/ESP-ONT-02/EXECUCAO\_LOOP.md}.
  \item \textbf{Linha de base técnica:} \texttt{docs/relatorios/ontologia/relatorio.pdf} (RT-ONT-01).
  \item \textbf{Taxonomia de base:} Alanoca, S., Gur-Arieh, S., Zick, T., Klyman, K. (2025). \emph{Comparing Apples to Oranges: A Taxonomy for Navigating the Global Landscape of AI Regulation}. FAccT~'25. doi:10.1145/3715275.3732059.
\end{itemize}

\end{document}
